%%==========================================================================
%%  Teaching Statement — Zheng Ma
%%==========================================================================
\documentclass[12pt, a4paper]{article}

%--- Line spacing -----------------------------------------------------------
\usepackage{setspace}
\setstretch{1.15}

%--- Encoding & fonts (LuaLaTeX) --------------------------------------------
\usepackage{fontspec}
\setmainfont{TeX Gyre Termes}
\setsansfont{TeX Gyre Heros}
\usepackage{microtype}

%--- CJK support (LuaLaTeX) -------------------------------------------------
\usepackage{luatexja-fontspec}
\setmainjfont{Noto Serif CJK SC}
\setsansjfont{Noto Sans CJK SC}

%--- Page geometry ----------------------------------------------------------
\usepackage[
  top=2.4cm, bottom=2.4cm,
  left=2.5cm, right=2.5cm,
  headheight=14pt
]{geometry}

%--- Color ------------------------------------------------------------------
\usepackage{xcolor}
\definecolor{accentcolor}{HTML}{800000}
\definecolor{linkcolor}{HTML}{1A3A5C}
\definecolor{sectionbg}{HTML}{F8F9FA}

%--- Math -------------------------------------------------------------------
\usepackage{amsmath, amssymb, amsthm, bm}
\usepackage{unicode-math}
\setmathfont{TeX Gyre Termes Math}

%--- Tables & Drawing -------------------------------------------------------
\usepackage{booktabs}
\usepackage{colortbl}
\usepackage{tabularx}
\usepackage{array}
\usepackage{float}
\usepackage{longtable}

%--- Headers / footers ------------------------------------------------------
\usepackage{fancyhdr}
\pagestyle{fancy}
\fancyhf{}
\renewcommand{\headrulewidth}{0.4pt}
\fancyhead[L]{\small\textit{Teaching Statement \,---\, Zheng Ma}}
\fancyhead[R]{\small\textit{April 2026}}
\fancyfoot[C]{\small\thepage}

%--- Section titles ---------------------------------------------------------
\usepackage{titlesec}
\titleformat{\section}
{\large\sffamily\bfseries\color{accentcolor}}{\thesection\enspace}{0em}{}
\titlespacing*{\section}{0pt}{3.2ex plus 0.8ex minus 0.4ex}{0.5ex plus 0.1ex}

\titleformat{\subsection}
{\normalsize\sffamily\bfseries\color{black!90}}{\thesubsection\enspace}{0em}{}
\titlespacing*{\subsection}{0pt}{2.2ex plus 0.5ex minus 0.2ex}{0.8ex}

%--- Lists ------------------------------------------------------------------
\usepackage{enumitem}
\setlist[itemize]{leftmargin=2em, itemsep=0.3em, topsep=0.5em, parsep=0.1em}
\setlist[enumerate]{leftmargin=2em, itemsep=0.3em, topsep=0.5em, parsep=0.1em}

%--- Hyperref ---------------------------------------------------------------
\usepackage[
  colorlinks=true,
  linkcolor=linkcolor,
  citecolor=linkcolor,
  urlcolor=linkcolor,
  pdftitle={Teaching Statement --- Zheng Ma},
  pdfauthor={Zheng Ma}
]{hyperref}

%--- Typography -------------------------------------------------------------
\usepackage{parskip}
\setlength{\parskip}{0.6em plus 0.2em}
\setlength{\parindent}{0pt}
\frenchspacing
\usepackage{xurl}

%%==========================================================================
%%  DOCUMENT
%%==========================================================================
\begin{document}

%%--------------------------------------------------------------------------
%%  TITLE BLOCK
%%--------------------------------------------------------------------------
\begin{center}
  {\LARGE\bfseries Teaching Statement}\\[0.8em]
  {\large Zheng Ma}\\[0.3em]
  {\normalsize
  Tenure-Track Associate Professor of Mathematics, Shanghai Jiao Tong University}\\[0.25em]
  {\normalsize \href{mailto:mazheng@sjtu.edu.cn}{mazheng@sjtu.edu.cn}}
\end{center}

\vspace{0.5em}
\noindent\rule{\textwidth}{0.4pt}
\vspace{0.2em}

%%--------------------------------------------------------------------------
%%  SECTION 1 — Teaching Philosophy
%%--------------------------------------------------------------------------
\section{Teaching Philosophy}

My approach to teaching mathematics and computation rests on a conviction that
has grown stronger with each passing year: the most lasting gift a teacher can
give is not a body of knowledge, but a \emph{way of thinking}.
Content becomes outdated; methods endure.
Every course I teach --- from undergraduate numerical analysis to graduate
high-performance scientific computing --- is organized around this principle.

\subsection{Teach the Method, Not the Content}

In a rapidly evolving technological landscape, specific algorithms or
numerical schemes may be superseded by new tools within years of a student's
graduation.
What cannot be superseded is the ability to \emph{ask the right question},
to \emph{break a hard problem into tractable pieces}, and to
\emph{reason about approximation and error}.
Rather than covering the maximum possible material, I deliberately slow down
at foundational ideas --- convergence, stability, well-posedness, the interplay
between continuous and discrete mathematics --- and ask students to derive,
question, and reconstruct.
In my \emph{Scientific Computing} and \emph{Numerical Analysis} courses,
for instance, students are required to implement algorithms from scratch
before they are permitted to use library routines.
Only by feeling the difficulty of the implementation do they truly internalize why
a method works, and where it might fail.

\subsection{Bidirectional Learning: Teaching as Conversation}

I regard every class as an opportunity to learn, not merely to instruct.
Students at Shanghai Jiao Tong University come from diverse backgrounds in
mathematics, engineering, and computer science; their questions and
misconceptions constantly reveal blind spots in my own understanding and
expose gaps in my explanations.
I welcome challenge, confusion, and pushback as productive signals.
Over the years, student questions have prompted me to revise derivations,
discover cleaner proofs, and occasionally reconsider aspects of my own
research.
Teaching, at its best, is a dialogue: the teacher shapes the student's
thinking, and the student sharpens the teacher's.

\subsection{Connecting Mathematics to Applications: Emphasizing \emph{Hands-On} Competence}

Pure mathematical theory, however beautiful, becomes inert without contact
with real problems.
I insist on grounding every abstract concept in a concrete application.
In \emph{Scientific Computing}, classical topics such as iterative solvers and
spectral methods are immediately paired with physical simulation exercises ---
heat diffusion, fluid flow, wave propagation.
In \emph{Hands-On Deep Learning (Python)}, students do not merely read about
neural networks; they build, train, and diagnose them on genuine datasets.
The course title captures the philosophy perfectly: ``\emph{动手}'' (hands-on)
is not an optional enrichment but the primary mode of learning.

I believe that mathematical education in the service of engineering and science
must produce graduates who can \emph{close the loop} between theorem and code,
between analysis and experiment.
A student who can prove a convergence theorem but cannot implement the method,
or who can train a model but cannot explain why it converges, has only half
the competence the modern world demands.

\subsection{Teaching for the AI Era: Breadth, Agility, and Judgment}

Artificial intelligence is reshaping what it means to be a scientist or
engineer.
Rote memorization and narrow specialization are becoming less valuable;
the ability to \emph{navigate unfamiliar territory quickly} is becoming
essential.
My teaching therefore emphasizes three meta-skills:

\begin{enumerate}
  \item \textbf{Learning how to learn.}
    I dedicate time in every course to discussing how to read technical
    literature efficiently, how to identify key ideas in a new paper or
    software library, and how to use AI tools as \emph{accelerators}
    rather than crutches.
    A student who can orient themselves in an unfamiliar field within days ---
    rather than months --- will always outpace one who knows a single
    domain deeply but cannot venture beyond it.

  \item \textbf{Prioritizing what matters.}
    Not all knowledge is equally important.
    I teach students to distinguish \emph{foundations} (which generalize
    across fields) from \emph{techniques} (which are domain-specific and
    often searchable) from \emph{trivia} (which need not be memorized at all).
    In an era when large language models can retrieve facts on demand,
    the scarce resource is \emph{judgment}: knowing what to ask, what to trust,
    and what to verify.

  \item \textbf{Width before depth.}
    My guiding heuristic for a student's early career is: \emph{go wider
    first, deeper with AI when you need it}.
    Broad exposure to mathematics, computation, physics, and engineering creates
    the cross-domain intuitions that spark creative research.
    Depth can be acquired on demand, especially with the assistance of modern
    AI tools, but breadth --- the ability to see unexpected connections ---
    must be cultivated progressively over years.
    I encourage my students to take courses outside their immediate specialty,
    to read across disciplines, and to resist the temptation of premature
    specialization.
\end{enumerate}

%%--------------------------------------------------------------------------
%%  SECTION 2 — Courses Taught
%%--------------------------------------------------------------------------
\section{Teaching Experience}

Since joining Shanghai Jiao Tong University in 2020, I have taught a range of
undergraduate and graduate courses spanning computational mathematics, numerical
analysis, and artificial intelligence.
The table below summarizes the courses to which I have made substantial
independent contributions.

\vspace{0.5em}

{\small
  \begin{longtable}{@{}p{5.8cm} l l r r@{}}
    \toprule
    \textbf{Course} & \textbf{Level} & \textbf{Term} & \textbf{Credits} & \textbf{Hours (personal)} \\
    \midrule
    \endhead
    \bottomrule
    \endfoot
    数值分析 (Numerical Analysis)
    & UG & 2025--26 S2 & 3 & 48 \\
    人工智能基础（A）(AI Fundamentals)
    & UG & 2025--26 S2 & 2 & 32 \\
    动手学深度学习（Python）
    & UG & 2025--26 S1 & 2 & 32 \\
    计算方法 (Computational Methods)
    & Grad & 2025--26 S1 & 3 & 48 \\[2pt]
    科学计算 (Scientific Computing)
    & UG & 2024--25 S2 & 3 & 48 \\
    研讨课~II (Seminar II)
    & UG & 2024--25 S2 & 3 & 48 \\
    人工智能导论 (Introduction to AI)
    & UG & 2024--25 S2 & 1 & 16 \\
    计算方法 (Computational Methods)
    & Grad & 2024--25 S1 & 3 & 48 \\
    微分方程的高性能计算
    & Grad & 2024--25 S1 & 2 & 32 \\[2pt]
    科学计算 (Scientific Computing)
    & UG & 2023--24 S2 & 3 & 48 \\
    科学问题的可计算建模
    & Grad & 2023--24 S2 & 1 & 16 \\
    计算方法 (Computational Methods)
    & Grad & 2023--24 S1 & 3 & 48 \\
    微分方程的高性能计算
    & Grad & 2023--24 S1 & 1 & 16 \\[2pt]
    科学计算 (Scientific Computing)
    & UG & 2022--23 S2 & 3 & 48 \\
    计算方法 (Computational Methods)
    & Grad & 2022--23 S1 & 3 & 48 \\[2pt]
    科学计算 (Scientific Computing)
    & UG & 2021--22 S2 & 3 & 48 \\
    计算方法 (Computational Methods)
    & Grad & 2021--22 S1 & 3 & 48 \\[2pt]
    微分方程数值解法 (Numerical PDEs)
    & Grad & 2020--21 S2 & 3 & 48 \\
    科学计算 (Scientific Computing)
    & UG & 2020--21 S2 & 3 & 48 \\
  \end{longtable}
}

Graduate sections of \emph{Computational Methods} (MATH6004) are delivered in
parallel Chinese and English to accommodate international students.
Undergraduate courses are all conducted in Chinese.

%%--------------------------------------------------------------------------
%%  SECTION 3 — Course Development
%%--------------------------------------------------------------------------
\section{Course Development and Innovation}

\subsection{Hands-On Deep Learning (Python)}

In 2025 I introduced \emph{Hands-On Deep Learning (Python)} (MATH2806), a
new required course for undergraduate mathematics majors.
The course is modeled on the open-source textbook \emph{Dive into Deep Learning}
(\url{https://d2l.ai}), whose central pedagogical idea --- every concept is
immediately paired with executable code --- aligns perfectly with my
philosophy of hands-on learning.
Students implement, from scratch, key architectures including multilayer
perceptrons, convolutional neural networks, recurrent networks, and
transformers, using PyTorch.
The course bridges the department's traditional numerical analysis curriculum
and the rapidly evolving practice of machine learning, and it reflects my
belief that a modern applied mathematician must be fluent in both worlds.

\subsection{Scientific Computing as a Bridge Course}

\emph{Scientific Computing} (MATH2802) has been my primary undergraduate
teaching responsibility since my arrival at SJTU.
Over successive offerings I have steadily increased the programming component,
migrating from MATLAB-based exercises to Python/NumPy/SciPy workflows.
Problem sets now include open-ended simulation projects --- modeling
epidemiological dynamics, simulating partial differential equations on irregular
domains, and benchmarking iterative solvers on sparse systems drawn from
finite element discretizations.
These projects require students to integrate mathematical analysis, algorithmic
thinking, and software engineering, and they have produced some of the most
creative and rigorous student work I have encountered.

\subsection{AI Literacy for Mathematics Students}

As AI tools become ubiquitous, I have incorporated structured discussions of
AI-assisted research and learning into multiple courses.
In \emph{Introduction to AI} (MATH1804) and the general-education course
\emph{AI Fundamentals} (MATH1116), I go beyond tool demonstrations to address
epistemic questions: How do you know when to trust an AI-generated result? How
do you efficiently decompose a learning task so that AI tools accelerate rather
than substitute for understanding?
These conversations are not asides; they are central to preparing students for
careers in which human--AI collaboration will be the norm.

%%--------------------------------------------------------------------------
%%  SECTION 4 — Future Directions
%%--------------------------------------------------------------------------
\section{Future Directions}

Looking ahead, I plan to develop a new graduate course on
\emph{Scientific Machine Learning} that unifies classical numerical analysis
with modern deep learning methodology --- covering physics-informed neural
networks, neural operators, and the mathematical foundations of generative
models, with a strong emphasis on provable guarantees and
hands-on computational projects.

I also intend to expand the use of project-based assessment across my courses,
replacing some conventional examinations with open-ended computational
investigations that require students to formulate, solve, and critically
evaluate a mathematical problem of their choosing.
This approach rewards precisely the skills I most want to cultivate:
breadth of vision, agility in learning, and rigorous judgment.

Finally, I am committed to making high-quality mathematical education in
scientific computing accessible beyond SJTU's campus.
Several of my courses already have openly available lecture notes and code
repositories, and I plan to continue developing these resources in collaboration
with colleagues and students.

%%--------------------------------------------------------------------------
%%  CLOSING
%%--------------------------------------------------------------------------
\vspace{1.5em}
\noindent\rule{\textwidth}{0.4pt}
\vspace{0.3em}

\noindent
{\small
  Teaching, for me, is inseparable from research: both are acts of inquiry,
  and both are most productive when conducted in a spirit of genuine curiosity
  and mutual respect.
  I am grateful every semester for the students who challenge my thinking,
  correct my mistakes, and remind me why mathematics matters.
}

\end{document}
